\documentclass{beamer}
\usepackage[utf8]{inputenc}
\usepackage[croatian]{babel}
\usepackage{amsmath}
\usepackage{amsfonts}
\usepackage{graphicx}
\usepackage{listings}
\usepackage{xcolor}

\definecolor{codegreen}{rgb}{0,0.6,0}
\definecolor{codegray}{rgb}{0.5,0.5,0.5}
\definecolor{codepurple}{rgb}{0.58,0,0.82}
\definecolor{backcolour}{rgb}{0.95,0.95,0.92}

\lstdefinestyle{mystyle}{
    backgroundcolor=\color{backcolour},   
    commentstyle=\color{codegreen},
    keywordstyle=\color{magenta},
    numberstyle=\tiny\color{codegray},
    stringstyle=\color{codepurple},
    basicstyle=\ttfamily\footnotesize,
    breakatwhitespace=false,         
    breaklines=true,                 
    captionpos=b,                    
    keepspaces=true,                 
    numbers=left,                    
    numbersep=5pt,                  
    showspaces=false,                
    showstringspaces=false,
    showtabs=false,                  
    tabsize=2
}
\lstset{style=mystyle}

\title{Prüferovi kodovi i problem "Clues"}
\author{Analiza teorije i implementacija rješenja}
\date{\today}

\begin{document}

\frame{\titlepage}

\begin{frame}
\frametitle{Uvod: Prüferovi kodovi}
\begin{itemize}
    \item \textbf{Prüferov niz (kod)} je jedinstveni niz koji predstavlja označeno stablo.
    \item Za stablo s $n$ čvorova označenih od $1$ do $n$, Prüferov kod ima duljinu $n-2$.
    \item Predstavlja \textbf{bijekciju} između skupa svih označenih stabala s $n$ čvorova i skupa svih nizova duljine $n-2$ nad abecedom {1, 2, ..., n}.
\end{itemize}
\end{frame}

\begin{frame}
\frametitle{Konstrukcija Prüferovog koda}
Algoritam za dobivanje Prüferovog koda iz stabla:
\begin{enumerate}
    \item Pronađi list (čvor stupnja 1) s najmanjom oznakom.
    \item Zabilježi oznaku njegovog jedinog susjeda u niz.
    \item Ukloni taj list i pripadajući brid iz stabla.
    \item Ponavljaj postupak dok u stablu ne ostanu točno 2 čvora.
\end{enumerate}
\vspace{0.5cm}
Ovo generira niz duljine $n-2$. Svaki čvor $v$ se u nizu pojavljuje točno $d(v) - 1$ puta, gdje je $d(v)$ stupanj čvora $v$ u stablu.
\end{frame}

\begin{frame}
\frametitle{Rekonstrukcija stabla iz Prüferovog koda}
Algoritam za dobivanje stabla iz Prüferovog koda $P$ duljine $n-2$:
\begin{enumerate}
    \item Generiraj skup $S$ koji sadrži sve moguće oznake čvorova $\{1, 2, \dots, n\}$.
    \item Ponavljaj sljedeće dok niz $P$ nije prazan:
    \begin{itemize}
        \item Neka je $v$ prvi element u nizu $P$.
        \item Pronađi najmanji element $u$ iz skupa $S$ koji se \textbf{ne nalazi} u nizu $P$.
        \item Dodaj brid $(u, v)$ u stablo.
        \item Ukloni $v$ s početka niza $P$ i ukloni $u$ iz skupa $S$.
    \end{itemize}
    \item Na kraju, u skupu $S$ ostat će točno 2 elementa (npr. $u$ i $v$). Dodaj zadnji brid $(u, v)$ u stablo.
\end{enumerate}
\vspace{0.3cm}
Ovaj algoritam jasno pokazuje da za svaki ispravan Prüferov niz duljine $n-2$ postoji jedinstveno stablo, čime je dokazana bijekcija.
\end{frame}

\begin{frame}
\frametitle{Cayleyjeva formula}
Kao izravna posljedica bijekcije Prüferovih kodova proizlazi \textbf{Cayleyjeva formula}.
\vspace{0.5cm}
\begin{block}{Teorem (Cayleyjeva formula)}
Broj različitih označenih stabala na $n$ čvorova je jednak $n^{n-2}$.
\end{block}
\vspace{0.5cm}
Zašto? Postoji $n^{n-2}$ različitih nizova duljine $n-2$ čiji su elementi iz skupa od $n$ mogućih vrijednosti. Zbog bijekcije, toliko ima i označenih stabala.
\end{frame}

\begin{frame}
\frametitle{Problem "Clues" (Codeforces 156D)}
\textbf{Opis problema:}
\begin{itemize}
    \item Zadan je graf s $n$ čvorova (tragovi) i $m$ bridova (veze).
    \item Graf se sastoji od $k$ povezanih komponenti.
    \item Potrebno je dodati \textbf{minimalan broj bridova} kako bi graf postao povezan (kako bi postojao put između svaka dva traga).
    \item Cilj: Izračunati \textbf{broj različitih načina} za dodavanje tih bridova, modulo dani broj. Dva načina su različita ako postoji brid koji je u jednom dodan, a u drugom nije.
\end{itemize}
\end{frame}

\begin{frame}
\frametitle{Analiza problema}
\begin{itemize}
    \item Minimalni broj bridova za povezivanje $k$ komponenti je točno $k - 1$.
    \item Dodavanjem $k-1$ bridova formiramo razapinjuće stablo (spanning tree) čiji su čvorovi zapravo originalne povezane komponente.
    \item Neka su veličine (broj čvorova) tih $k$ povezanih komponenti $s_1, s_2, \dots, s_k$.
    \item Spajanjem komponente $i$ s komponentom $j$, biramo jedan čvor iz prve komponente i jedan iz druge. To možemo napraviti na $s_i \times s_j$ načina.
\end{itemize}
\end{frame}

\begin{frame}
\frametitle{Uopćena Cayleyjeva formula}
Broj načina za povezivanje $k$ komponenti u jedno stablo definiran je formulom:
\[ N = n^{k-2} \prod_{i=1}^k s_i \]
\vspace{0.2cm}
Ovo je poznato kao uopćena Cayleyjeva formula. Kako do nje dolazimo?
\end{frame}

\begin{frame}
\frametitle{Dokaz pomoću Prüferovih kodova}
\begin{itemize}
    \item Razmotrimo stablo gdje su čvorovi komponente. Neka komponenta $i$ ima stupanj $d_i$ u tom novom stablu.
    \item Broj načina da odaberemo konkretne vrhove iz komponenti za to specifično stablo je $\prod_{i=1}^k s_i^{d_i}$.
    \item Zbroj po svim mogućim stablima na $k$ čvorova:
    \[ \sum_{T} \prod_{i=1}^k s_i^{d_i} = \prod_{i=1}^k s_i \sum_{T} \prod_{i=1}^k s_i^{d_i - 1} \]
    \item U Prüferovom nizu duljine $k-2$, čvor $i$ se pojavljuje $d_i - 1$ puta. Sumiranje preko svih mogućih Prüferovih nizova je ekvivalentno:
    \[ \sum_{x_1, \dots, x_{k-2}} s_{x_1} s_{x_2} \dots s_{x_{k-2}} \]
\end{itemize}
\end{frame}

\begin{frame}
\frametitle{Završetak dokaza}
\begin{itemize}
    \item Gornja suma je po \textbf{multinomnom teoremu} jednaka:
    \[ (s_1 + s_2 + \dots + s_k)^{k-2} \]
    \item Budući da je suma veličina svih komponenti jednaka ukupnom broju čvorova $n$ ($s_1 + s_2 + \dots + s_k = n$), dobivamo:
    \[ n^{k-2} \]
    \item Konačan rezultat je dakle:
    \[ n^{k-2} \prod_{i=1}^k s_i \]
    \item Poseban slučaj: Ako je $k=1$, broj načina je $1 \pmod{\text{modulo}}$.
\end{itemize}
\end{frame}

\begin{frame}[fragile]
\frametitle{Rješenje: Implementacija u Rustu (1/2)}
\begin{lstlisting}[basicstyle=\ttfamily\tiny]
// Nalazenje povezanih komponenti pomocu BFS-a
let mut visited = vec![false; n];
let mut components = Vec::new();

for i in 0..n {
    if !visited[i] {
        let mut count = 0u64;
        let mut q = VecDeque::new();
        q.push_back(i);
        visited[i] = true;
        while let Some(u) = q.pop_front() {
            count += 1;
            for &v in &adj[u] {
                if !visited[v] {
                    visited[v] = true;
                    q.push_back(v);
                }
            }
        }
        components.push(count);
    }
}
\end{lstlisting}
\end{frame}

\begin{frame}[fragile]
\frametitle{Rješenje: Implementacija u Rustu (2/2)}
\begin{lstlisting}
let k = components.len() as u64;

// Poseban slucaj
if k == 1 {
    println!("{}", 1 % modulo);
    return;
}

// Izracun formule: n^(k-2) * product(s_i)
let mut ans = pow(Mod::new(n as u64, modulo), k - 2);
for &c in &components {
    ans = ans * Mod::new(c, modulo);
}

println!("{}", ans.value());
\end{lstlisting}
\end{frame}

\begin{frame}
\frametitle{Složenost}
\begin{itemize}
    \item \textbf{Vremenska složenost:} 
    \begin{itemize}
        \item Učitavanje grafa i pronalazak komponenti (BFS): $\mathcal{O}(n + m)$.
        \item Modularno potenciranje $n^{k-2}$: $\mathcal{O}(\log k)$.
        \item Množenje veličina komponenti: $\mathcal{O}(k)$.
        \item Ukupna vremenska složenost: $\mathcal{O}(n + m + \log k)$.
    \end{itemize}
    \item \textbf{Prostorna složenost:} 
    \begin{itemize}
        \item Predstavljanje grafa (liste susjedstva): $\mathcal{O}(n + m)$.
        \item Pomoćni nizovi (\texttt{visited}, \texttt{components}, BFS red): $\mathcal{O}(n)$.
        \item Ukupna prostorna složenost: $\mathcal{O}(n + m)$.
    \end{itemize}
\end{itemize}
\end{frame}

\begin{frame}
\begin{center}
    \Huge Pitanja?
\end{center}
\end{frame}

\end{document}
